% Generated draft: requires mathematical and source review.
\documentclass[11pt]{article}
\usepackage[margin=25mm]{geometry}
\usepackage{fontspec}
\setmainfont{Latin Modern Roman}
\usepackage{amsmath,amssymb,mathtools,unicode-math}
\setmathfont{Latin Modern Math}
\usepackage{enumitem,needspace,xurl,hyperref}
\setlength{\parindent}{0pt}
\setlength{\parskip}{5pt}
\setlength{\emergencystretch}{3em}
\newcommand{\N}{\mathbb{N}}
\newcommand{\Z}{\mathbb{Z}}
\newcommand{\Q}{\mathbb{Q}}
\newcommand{\R}{\mathbb{R}}
\newcommand{\C}{\mathbb{C}}
\newcommand{\F}{\mathbb{F}}
\newcommand{\A}{\mathbb{A}}
\newcommand{\PP}{\mathbb{P}}
\newcommand{\E}{\mathbb{E}}
\newcommand{\eps}{\varepsilon}
\newcommand{\dd}{\,\mathrm{d}}
\begin{document}
\section*{UnsolvedMath ID: 800355}
\textbf{Source problem code:} catalogue entry
\subsection{Definitions and mathematical statement}
Let \(D=\{n\in\Z:n\geq 5\}\), equipped with the topology induced by the \(2\)-adic metric. Write formally
\[
a(n)=T_{113}(n,2n-4),\qquad n\in D.
\]
The supplied record does not define \(T_{113}\). In particular, it does not specify the objects counted, the meaning of the subscript \(113\), the parameters, or the boundary conditions. Thus this notation cannot be given a complete combinatorial definition from the supplied evidence alone.

Assuming that the intended quantities \(a(n)\) are integers, as the source's congruence observations presuppose, the question is whether \(a:D\to\Z\) is continuous for the induced \(2\)-adic topologies. Explicitly, ordinary continuity on this domain asks whether
\[
\forall n\in D\;\forall r\in\Z_{\geq 1}\;\exists s\in\Z_{\geq 0}\;\forall m\in D,\qquad
m\equiv n\pmod{2^s}\ \Longrightarrow\ a(m)\equiv a(n)\pmod{2^r}.
\]
Here \(s\) may depend on both \(n\) and \(r\). The source does not specify whether it instead intends uniform continuity, in which \(s\) must depend only on \(r\), or a continuous extension to the \(2\)-adic integers. These stronger interpretations are not imposed in the displayed formulation.

The supplied comment reports apparent constancy modulo \(2\) and modulo \(4\), apparent \(2\)-periodicity modulo \(8\), and apparent \(8\)-periodicity modulo \(16\), with repeating pattern
\[
4,\ 4,\ 4,\ 4,\ 4,\ 12,\ 4,\ 12.
\]
These are reported observations, not established hypotheses or conclusions. The record does not explicitly identify the starting index of the displayed pattern. It states that \(n=4\) breaks the pattern; that index is excluded from the question. The registry labels the problem open, but that status has not been independently verified.
\subsection{Short English statement}
Propp asks whether the sequence \(T_{113}(n,2n-4)\), for integers \(n\geq 5\), is continuous in the \(2\)-adic topology. The congruence formulation above uses ordinary continuity. A complete statement still requires the missing definition of \(T_{113}\) and confirmation of the intended continuity convention. The reported modular patterns are exploratory observations.
\subsection{Sources}
\begin{itemize}
\item Local UnsolvedMath registry, permanent ID \(800355\), exact source problem code \texttt{catalogue entry}, title ``Trimer Covers in the Triangular Grid --- Problem 15.'' The supplied record attributes the problem to Propp (2022) and labels it open. Its recorded upstream revision is \texttt{d371ee917ac9647fbb06cd6fc529deec6f3dda82}.
\item Propp, \emph{Trimer Covers in the Triangular Grid: Twenty Mostly Open Problems} (2022), source item ``Problem 15 (\#7)'': \url{https://arxiv.org/abs/2206.06472}. The supplied provenance records access on 2026-10-16 and extraction by \texttt{arxiv\_latex}; these are supplied provenance details, not a claim of independent access or verification.
\item OEIS entry A364416, cited in the supplied problem comment: \url{https://oeis.org/A364416}. Its contents have not been independently inspected for this draft.
\end{itemize}
\end{document}
